\mode*
\section{Errors and debugging the document}
\label{sec:errors}

The second educational objective is error handling: reading TeX's error
messages, locating the offending source, and the recompile loop as a
debugging loop.

\subsection{Documented difficulties}

The systematic round (\cref{app:protocol}) again found no LaTeX-specific
studies of error handling by learners, but the adjacent literature on
\emph{programming} error messages is well developed, and its findings are
suggestive for TeX errors.
Interpreting error messages is crucial for making progress, yet the
messages are often cryptic and pose a barrier to novices, who repeat the
same errors \autocite{Becker2016CEM,Traver2010CompilerErrors}, and poorly
designed messages hurt novices more than experts
\autocite{Traver2010CompilerErrors}; their readability turns on such
factors as message length and jargon \autocite{Denny2021ErrorReadability}.
TeX's terse, idiosyncratic error reports (\enquote{\texttt{Undefined
control sequence}}, line numbers pointing past the real fault) are at
least as cryptic, and we conjecture the same barrier for learners.
Most directly for the aspects below, novices do not reliably act on the
\emph{first} error: in one activity-log study students addressed the first
error only about half the time, and doing so correlated with better grades
\autocite{Munson2016CompilerErrors}; and students who do read the message
generally make effective changes afterwards
\autocite{Prather2017Interaction}---so reading the message, and reading the
first one, are the moves to teach.
The markup-education literature shows the other side of the same coin:
where the system does \emph{not} report errors, they persist.
Browsers render HTML and CSS however broken, so in a lab study
\enquote{the errors remain latent and unresolved, reinforcing faulty
understandings} \autocite{Park2013ErrorTaxonomy}\footnote{Quoted from
\textcite[pp.~107--108]{Park2014openHTML}, which reports the same
study.}, and in a course-wide log study nearly all students' syntax errors
stayed unresolved for weeks, while \enquote{nearly all of the errors that
were detected through validation were eventually corrected}
\autocite{Park2015SyntaxErrors}.
TeX stops and reports; the report is the resource, if it is read.
% TODO(#2): Course data: which TeX errors stop datintro students, and what do
% they do next?

\subsection{Candidate critical aspects}

% XXX Preliminary.
\begin{itemize}
  \item \emph{An error message is information}: it names a place and a
    cause, in that order of reliability.
  \item \emph{The first error counts}: later errors are often
    consequences.
  \item \emph{Errors are normal}: the recompile loop is the working mode,
    not a failure mode. (Debugging as learning: cf.\ the companion paper
    \autocite{VTDebug}.)
\end{itemize}

\subsection{Patterns of variation}

\begin{itemize}
  \item Contrast (one defect varies): plant a single defect (missing
    brace, misspelled command, missing package) in an otherwise working
    document; read the message each produces.
  \item Generalisation (defect invariant, document varies): the same
    missing-brace error in different documents---the message's shape is
    the invariant to discern.
  \item Fusion: two planted defects; find and fix both using
    first-error-first.
\end{itemize}
